\section{相关工作}

\subsection{受限理性元推理与智能体过度思考}
机载计算资源的分配问题源于人工智能领域的受限理性与元推理理论，其核心在于评估“计算的价值”是否超过计算所耗费的时间与能量代价\cite{horvitz1987reasoning,russell1991right}。在机器人运动规划与实时控制中，研究者通过学习停止规则或自适应计算预算，决定何时终止路径搜索并执行当前最优动作，或将规划与移动并行交织\cite{xiao2020thinking,sung2021quit,elboher2023formal,magnino2026metareasoning}。然而，当内部环境模型本身存在严重偏差时，延长规划视野或增加搜索深度反而会放大模型误差，诱发决策分析中的“优化器诅咒”，使规划器陷入无益的过度思考\cite{smith2006curse,jiang2015horizon,cuadron2025overthinking}。近期在智能体推理与任务消歧中也发现了高度一致的规律：当外部环境状态未知或指令本身存在多重解释时，增加测试时计算无法替代与环境的交互或主动向用户寻求澄清\cite{hunt2021formalizing,shen2025thinking,kobalczyk2025active}。与上述仅关注何时停止搜索的元推理方法不同，本文将模糊性明确拆解为求解器粗糙引起的数值歧义与地图失真引起的环境歧义，从而在继续计算、获取观测与保守行动之间建立定量切换准则。

\subsection{面向决策的信息获取与多解性判别}
并非所有减少环境不确定性的观测都能改善特定决策，唯有能够区分竞争假设的信息才具备决策价值\cite{chernoff1959sequential,box1967discrimination}。在统计实验设计与多臂赌博机理论中，最优臂识别与信息导向采样通过度量备选动作之间的置信重叠或信息比，专门收集能够尽快排除次优选项的样本，而非均匀降低全局参数方差\cite{garivier2016optimal,russo2018thompson,russo2018ids}。在大规模贝叶斯反问题中，目标导向实验设计同样强调只针对感兴趣的目标泛函优化传感器布放\cite{attia2018goal}。与之呼应，信息间隙决策理论与机器学习中的预测多解性研究指出，当多个与训练数据同样相容的模型对同一样本给出截然相反的预测时，单靠点估计无法保证决策安全，必须考察给定不确定性半径内的最坏表现\cite{benhaim2006infogap,marx2020predictive}。不同于在离散假设空间或线性反问题中计算期望信息增益的传统策略，本文利用包络定理直接在连续无散海流场上沿梯度搜寻翻转二元航路选择的物理反例，以最坏后悔直接度量环境歧义的严重程度。

\subsection{不确定海洋流场中的航路规划与自适应采样}
强海流环境下的自主水下航行器路径规划长期受到海洋工程界关注，经典方法利用图搜索或水平集方程在已知确定性流场中求取时间最优轨迹\cite{dijkstra1959note,lolla2014levelset}。为应对海洋预报的不确定性，研究者相继提出了基于集合预报的期望与风险折中路径规划、动态正交方程驱动的随机时间最优控制，以及显式纳入预报误差协方差的鲁棒路径搜索\cite{wang2016path,subramani2018stochastic,kularatne2018forecast}。与此同时，海洋自适应采样研究利用高斯过程或集合卡尔曼滤波在线更新叶绿素或温盐流场估计，引导航行器前往信息增益高的特征区域\cite{anderson2001eakf,srinivas2012gpucb,hollinger2014sampling,fossum2019}。然而，现有不确定流场规划大多将集合样本或概率分布作为固定输入来寻找单一路径，既不评估有限集合是否漏掉了翻转拓扑走向的极端流场，也不回答在出发前或航行中是否值得付出时间代价去换取新海流图；本文则将多年海洋再分析大数据转化为合理流场集合的物理边界，专门解决粗分辨率海流图下“算、测、走”的闭环元决策问题。
